\documentclass[10pt,a4paper]{amsart}
\usepackage[margin=19mm]{geometry}
\usepackage{amsmath,amssymb,mathtools}
\usepackage[scr=rsfs,cal=euler]{mathalfa}
\usepackage{xfrac}
\usepackage[unicode,hidelinks,bookmarks=false]{hyperref}
\newcommand{\ie}{i.e.\ }
\newcommand{\cf}{cf.\ }
\newcommand{\resp}{respectively\ }
\newcommand{\Subsection}{\subsection}
\providecommand{\nobreakdash}{\nobreak\hbox{-}}
\setlength{\emergencystretch}{2em}
\multlinegap0pt
\usepackage{amssymb}
\usepackage{mathtools}
\usepackage[shortlabels]{enumitem}
\usepackage{float}
\usepackage{tikz}
\usepackage{xcolor}
\newtheorem{theorem}{Theorem}
\newtheorem{corollary}{Corollary}
\newtheorem{claim}{Claim}
\usepackage{cleveref}
\crefname{theorem}{Theorem}{Theorems}
\crefname{corollary}{Corollary}{Corollarys}
\crefname{claim}{Claim}{Claims}
\newcommand{\R}{\mathbb R}
\newcommand{\ang}[1]{\langle#1\rangle}
\newcommand{\bin}{\{0,1\}}
\newcommand{\cH}{\mathcal H}
\newcommand{\cS}{\mathcal S}
\newcommand{\floor}[1]{\lfloor#1\rfloor}
\renewcommand{\ge}{\geqslant}
\renewcommand{\le}{\leqslant}
\title{Nondegenerate hyperplane covers of the hypercube}
\author{Lisa Sauermann, Zixuan Xu}
\date{Source: arXiv:2507.00773v2 (CC BY 4.0)}
\begin{document}
\maketitle

\noindent\emph{Source and licence.} Nondegenerate hyperplane covers of the hypercube. Version arXiv:2507.00773v2 (CC BY 4.0), distributed by arXiv under the Creative Commons Attribution 4.0 International licence, http://creativecommons.org/licenses/by/4.0/, as recorded in the arXiv licence metadata for this exact version and shown on the abstract page https://arxiv.org/abs/2507.00773v2. Full licence notice: \texttt{licenses/arxiv-2507.00773-CC-BY-4.0.txt}.\par\smallskip

\noindent\textit{Editorial scope.} Counted unit: the article's corollary cor:slicing (source0.tex lines 233--235). The article's complete original proof of it is delivered with it (source0.tex lines 313--334, one \texttt{proof} environment of 22 lines, which the article places away from the statement). No mathematics is written, repaired or re-derived: the statement and the proof are the article's own lines, in the article's own notation, and no retained line is substituted or re-flowed.  Retained uncounted as the source-local context the proof needs: lines 212--214.  Nothing was omitted from any retained span, so the package carries no omission record.  No retained line contains a citation, so the wrapper carries no bibliography and no bibliography entry.  No retained line contains a figure environment, an image-inclusion command or drawing code, so no image is delivered and the package declares no asset dependency.  Editorial formatting only, in addition: an AMS \texttt{amsart} preamble that restates the article's own declarations and macros used by the retained lines, namely the environments \texttt{theorem}, \texttt{corollary}, \texttt{claim} on separate counters, together with the standard packages those lines need.  The retained environments print in this document as Theorem 1, Corollary 1, Claim 1 in order of appearance, whereas the article prints the same items with the numbering of its own full document.  The complete native source of the article as distributed in the arXiv e-print for this exact version is bundled unchanged as \texttt{sources/180837/source0.tex}.
\begin{theorem}\label{thm:main}
    Let $\cH$ be a collection of hyperplanes in $\R^n$ satisfying the following condition: for every $v\in \bin^n$ and every $i=1,\dots,n$, there exists a hyperplane $H\in \mathcal{H}$ through $v$ such that the $i$-th coordinate of the normal vector of $H$ is non-zero. Then we must have $|\cH|\ge n/2$.
\end{theorem}

\begin{corollary}\label{cor:slicing}
    Let $C$ be a positive integer and let $\cH $ be a collection of hyperplanes in $\mathbb{R}^n$ with normal vectors in $\{-C, \dots, C\}^n$, such that every edge of the $n$-dimensional hypercube $[0,1]^n$ is sliced by a hyperplane in $\cH$. Then $|\cH|\ge n/(4C)$.
\end{corollary}

\begin{proof}[Proof of \cref{cor:slicing}]
Let $\cH$ be a set of hyperplanes in $\mathbb{R}^n$ with normal vectors in $\{-C,\dots, C\}^n$ such that every edge of the $n$-dimensional hypercube is sliced by at least one of the hyperplanes in $\cH$. Consider the set of hyperplanes $\cH'$ in $\mathbb{R}^n$ constructed from $\cH$ as follows: For each hyperplane $H\in \cH$ with hyperplane equation $\ang{a,x} = b $ and normal vector $a\in \{-C,\dots, C\}^n$, let $\cS_H$ denote the set of the $2C$ hyperplanes with hyperplane equations $\ang{a,x} = \floor{b} + z$ for $z\in \{-(C-1),\dots, C\}$. Let $\cH' = \bigcup_{H\in \cH} \cS_H$, and note that then we have $|\cH'| \le 2C\cdot |\cH|$.



\begin{claim}\label{claim:reduction}
    $\cH'$ is a collection of hyperplanes satisfying the condition in \cref{thm:main}.
\end{claim}

\begin{proof}
    Let $v\in \bin^n$ be a vertex and let $i\in \{1,\dots,n\}$. There must be a hyperplane $H\in \cH$ slicing the edge adjacent to $v$ along the $i$-th coordinate direction. Let $v'$ be the other endpoint of this edge (then $v$ and $v'$ only differ in the $i$-th coordinate). Let $H$ be described by the equation $\ang{a,x} = b$  with $a\in \{-C,\dots, C\}^n$ and $b\in \R$. Since $H$ slices the edge $vv'$ of the hypercube, the vertices $v$ and $v'$ lie on different sides of the hyperplane $H$, so $\ang{a,v} - b$ and $\ang{a,v'} - b$ have different signs (meaning one of the numbers $\ang{a,v} - b$ and $\ang{a,v'} - b$ is positive, while the other one is negative). 
    On the other hand, $v$ and $v'$ differ only in the $i$-th coordinate (and the $i$-th coordinate of one of these vertices is $1$, while for the other one it is $0$), so we have $|\ang{a,v}-\ang{a,v'}|=|a_i|\le C$. Therefore we obtain
    \[|(\ang{a,v} - b) - (\ang{a,v'} - b)| = |\ang{a,v}-\ang{a,v'}| =|a_i| \le C,\]
    and since $\ang{a,v} - b$ and $\ang{a,v'} - b$ have different signs, we can conclude that $a_i\ne 0$ and $|\ang{a,v} - b|<C$ (indeed, if $|\ang{a,v} - b|\ge C$, then, due to the inequality above,  $\ang{a,v'} - b$ would need to be zero or have the same sign as $\ang{a,v} - b$). Thus, we have $\ang{a,v}\in (b-C,b+C)$, and since $a$ and $v$ are integer vectors (and $C$ is a positive integer), this implies
    \[\ang{a,v}\in (b-C,b+C)\cap \mathbb{Z}\subseteq \{\floor{b-C}+1,\dots, \floor{b+C}\}=\{\floor{b}-(C-1),\dots, \floor{b}+C\}.\]
    This means that $\ang{a,v}=\floor{b}+z$ for some $z\in \{-(C-1),\dots, C\}$, and so $v$ is contained in one of the hyperplanes in $\cS_H\subseteq \cH'$. Each of these hyperplanes has normal vector $a$, and the $i$-th coordinate of $a$ is $a_i\ne 0$. Thus, $v$ is indeed contained in a hyperplane in $\cH'$ whose normal vector has a non-zero entry in the $i$-th coordinate.
\end{proof}  



By \cref{thm:main}, we have $|\cH'| \ge n/2$. Since $|\cH'|\le 2C\cdot |\cH|$ by construction, we can conclude that $|\cH|\ge n/(4C)$, as desired.
\end{proof}

\end{document}
